% Original analytic teaching diagram, 2026-11-12.
% Sensor joint table (rows y, columns x; symmetric table of chapter 7.6) and the kernel
% K(y,{x}) obtained by dividing each row by its sum p_Y(y)=(0.3,0.4,0.3).
% The column x=2 of the kernel gives P(X=2|Y=y) = 0, 1/4, 2/3.
\definecolor{VB}{HTML}{7188A3}\definecolor{VBF}{HTML}{E5EAF0}
\definecolor{VG}{HTML}{829C8B}\definecolor{VGF}{HTML}{E8EEE9}
\definecolor{VY}{HTML}{B4A16F}\definecolor{VYF}{HTML}{F1ECDD}
\definecolor{VR}{HTML}{A7787D}\definecolor{VRF}{HTML}{F0E5E6}
\definecolor{VP}{HTML}{9588A5}\definecolor{VPF}{HTML}{ECE8F0}
\definecolor{VI}{HTML}{3B4252}\definecolor{VM}{HTML}{6B7280}\definecolor{VL}{HTML}{4C4D4F}
\begin{tikzpicture}[x=1mm,y=1mm,>=Stealth,font=\figurefont\fontsize{8}{10}\selectfont,
lab/.style={text=VM,font=\figurefont\fontsize{7.5}{9}\selectfont}]
\def\cw{9}\def\ch{7}
% ---------- (a) joint table ----------
\foreach \r/\a/\b/\c in {0/0.20/0.10/0,1/0.10/0.20/0.10,2/0/0.10/0.20}{
  \node[text=VI] at (0.5*\cw,-\r*\ch-0.5*\ch) {$\a$};
  \node[text=VI] at (1.5*\cw,-\r*\ch-0.5*\ch) {$\b$};
  \node[text=VI] at (2.5*\cw,-\r*\ch-0.5*\ch) {$\c$};
}
\draw[VL!55,line width=.4pt] (0,0) grid[xstep=\cw,ystep=\ch] (3*\cw,-3*\ch);
\draw[VL,line width=.7pt] (0,0) rectangle (3*\cw,-3*\ch);
\foreach \k in {0,1,2}{
  \node[text=VI] at (\k*\cw+0.5*\cw,3.5) {$x=\k$};
  \node[text=VI,left] at (-0.5,-\k*\ch-0.5*\ch) {$y=\k$};
}
\foreach \r/\s in {0/0.3,1/0.4,2/0.3}{
  \node[text=VI] at (3*\cw+7,-\r*\ch-0.5*\ch) {$\s$};
}
\node[text=VI] at (3*\cw+7,3.5) {$p_Y(y)$};
\draw[VL,line width=.6pt] (3*\cw+2.5,0) rectangle (3*\cw+11.5,-3*\ch);
\node[text=VI] at (1.5*\cw+4,-3*\ch-6) {(a)};
% ---------- arrow ----------
\draw[->,VL,line width=.7pt] (3*\cw+14,-1.5*\ch) -- (3*\cw+22,-1.5*\ch);
\node[lab,align=center] at (3*\cw+18,-1.5*\ch+4.5) {除以行和};
% ---------- (b) kernel ----------
\begin{scope}[xshift=55mm]
\fill[VBF] (2*\cw,0) rectangle (3*\cw,-3*\ch);
\foreach \r/\a/\b/\c in {0/\tfrac23/\tfrac13/0,1/\tfrac14/\tfrac12/\tfrac14,2/0/\tfrac13/\tfrac23}{
  \node[text=VI] at (0.5*\cw,-\r*\ch-0.5*\ch) {$\a$};
  \node[text=VI] at (1.5*\cw,-\r*\ch-0.5*\ch) {$\b$};
  \node[text=VI] at (2.5*\cw,-\r*\ch-0.5*\ch) {$\c$};
}
\draw[VL!55,line width=.4pt] (0,0) grid[xstep=\cw,ystep=\ch] (3*\cw,-3*\ch);
\draw[VL,line width=.7pt] (0,0) rectangle (3*\cw,-3*\ch);
\draw[VB,line width=.8pt] (2*\cw,0) rectangle (3*\cw,-3*\ch);
\foreach \k in {0,1,2}{
  \node[text=VI] at (\k*\cw+0.5*\cw,3.5) {$x=\k$};
}
\foreach \r in {0,1,2}{
  \node[lab] at (3*\cw+5,-\r*\ch-0.5*\ch) {和$=1$};
}
\node[text=VI] at (1.5*\cw,-3*\ch-6) {(b)};
\end{scope}
\end{tikzpicture}
